% TOP_PROBLEM: {"id": 10000011, "title": "Cheeger constants of nets in transitive graphs", "category_id": 3, "category": {"id": 3, "name": "graph_theory", "display_name": "Graph Theory", "description": "Problems involving graphs, networks, and their properties.", "slug": "graph-theory", "order_index": 3, "created_at": "2026-07-31T15:26:25.671Z"}, "status": "open", "problem_number": "AMR-099-0011", "set_id": 13, "statusQualification": "Uses the primary separation threshold k and adjacency threshold 2k, and states the finite-graph half-size convention explicitly."}
% FIRST_POSTED: 2026-10-01
% ENTRY_KIND: statement_only
% UnsolvedMath ID 10000011; source catalogue code AMR-099-0011
% !TeX program = lualatex
% Definition and sources only; this file is not an LLM solution attempt.
\documentclass[11pt,a4paper]{article}
\usepackage[margin=25mm]{geometry}
\usepackage{fontspec}
\setmainfont{lmroman10-regular.otf}[BoldFont=lmroman10-bold.otf,
 ItalicFont=lmroman10-italic.otf,BoldItalicFont=lmroman10-bolditalic.otf]
\usepackage{amsmath,amssymb,mathtools}
\usepackage{unicode-math}
\setmathfont{latinmodern-math.otf}
\AtBeginDocument{\let\setminus\smallsetminus}
\usepackage{xurl,hyperref}
\hypersetup{colorlinks=true,urlcolor=blue}
\setcounter{secnumdepth}{0}
\setlength{\parindent}{0pt}
\setlength{\parskip}{5pt}
\setlength{\emergencystretch}{3em}
\begin{document}
\section*{Cheeger constants of nets in transitive graphs}\label{10000011}
{\small UnsolvedMath ID 10000011 · AMR-099-0011 · Graph Theory · AMR Open Problem Lists}
\subsection{Definitions and mathematical statement}
For a graph $H$, write $\partial_H A$ for the vertices outside $A$ adjacent to some vertex of $A$. Define
\[
h(H)=\inf_A\frac{|\partial_H A|}{|A|},
\]
where $A$ ranges over nonempty finite vertex sets if $H$ is infinite, and over $0<|A|\le |V(H)|/2$ if $H$ is finite. Distances below are shortest-path distances with every edge of length one.

Let $G$ be a connected locally finite vertex-transitive graph and let $k\ge1$ be an integer. Choose any set $N\subseteq V(G)$ maximal under inclusion subject to $d_G(x,y)\ge k$ for distinct $x,y\in N$. The associated $k$-net graph $G_k$ has vertex set $N$, with an edge between two distinct members whenever their distance in $G$ is at most $2k$. Maximality implies every vertex of $G$ lies within distance less than $k$ of $N$.

Whenever $G_k$ has at least two vertices, must
\[
h(G_k)\ge h(G)?
\]
The assertion quantifies over every maximal separated set, not just the existence of a particularly well-chosen net. It covers both finite and infinite transitive graphs. Its comparison uses vertex expansion, so replacing the boundary by an edge count would be a different numerical assertion. The net retains adjacency at scale $2k$ while discarding most original vertices; the question is whether this coarsening can ever reduce the Cheeger constant in the transitive setting.
\subsection{Short English statement}
Can passage to any maximal graph net reduce the vertex Cheeger constant of a transitive graph?
\subsection{Sources}
\begin{itemize}
\item[S1] ULAM.AI and the UnsolvedMath Contributors, supplied problems.json, record 10000011 (AMR-099-0011), AMR Open Problem Lists. Catalogue identity and source transcription; CC BY 4.0.
\url{https://huggingface.co/datasets/ulamai/UnsolvedMath}.
\item[S2] Benjamini - Coarse Geometry and Randomness (Saint-Flour notes), Open problem 1.69, PDF page 15. Original source indexed by AMR-099-0011.
\url{https://www.wisdom.weizmann.ac.il/~itai/stflouraug24.pdf#page=15}.
\item[S3] I. Benjamini, Coarse Geometry and Randomness, primary Saint-Flour notes; the numbered question and its surrounding definitions.
\url{https://arquivo.pt/noFrame/replay/20201231041548id_/http://www.wisdom.weizmann.ac.il/~itai/stflouraug24.pdf}.
\end{itemize}
\end{document}
